\chapter{Kernel metadata and the image}\label{ch:kernel-metadata-image}

A compute kernel reaches the GPU as a binary archive whose machine code is accompanied by metadata
sections the driver reads when it builds a pipeline: which special registers the kernel reads, how its buffers are
bound, whether it uses the tensor unit, what its constant program writes, how much threadgroup memory it declares.
Apple documents none of these sections.

The GPU executes the project's machine code, and a change made to Metal's copy of it after pipeline creation reaches
execution. Pipeline creation validates only the container's format: undecodable code, a missing end of program and
impossible metadata values all build a pipeline (\cref{sec:object-library-archive}). Mutating one field at a time,
against a control that could fail, separated the fields the driver reads and acts on (the tensor-unit enable, the
system-register vector, the binding records) from those it reads only for their presence and those it ignores.
Damaged structure is not refused: it crashes the driver inside pipeline creation, and an earlier pointer-slot defect
hung the GPU. For these reasons the compiler computes metadata from the program it emitted instead of copying it from a
neighbouring kernel, and it refuses layouts it has not witnessed.

Names used for the image and its parts:

\begin{itemize}
\item \emph{Image}: the delivered \code{scan.arc.metallib}. \emph{Object}: the \code{applegpu} Mach-O inside it. \emph{Section}: a Mach-O section
  such as \code{\_\_GPU\_METADATA,\_\_compute}.
\item \emph{Per-kernel slot N}: field N of the per-kernel FlatBuffers table in \code{\_\_GPU\_METADATA} (\cref{sec:gpu-metadata-slots}). Not to be
  confused with scoreboard slots (\cref{sec:scoreboard}).
\item \emph{Constant program}: Apple's per-dispatch program \code{agc.main.constant\_program} (\cref{sec:constant-program-record}). \emph{Uniform} uK: the
  uniform register file (\cref{sec:register-classes}).
\item \emph{Kind-3 record}, \emph{kind-6 record}, \emph{argument map}: records of the constant-program description (\cref{sec:constant-program-record}).
\item \emph{Class}: a measured metadata layout (which vtable sits where, which slots exist) that this project can author
  (\cref{sec:abi-classes}).
\item \emph{Contract}: the compiler's statement of a program (\code{program.contract()}: code identity, entry, bindings, instruction
  inventory) and of its emission choices (\code{program.abi()}: register count, forms, system-register set, metadata class,
  launch) (MM~\mm{23.5}).
\end{itemize}

\section{\code{\_\_GPU\_METADATA} slots}\label{sec:gpu-metadata-slots}

\code{\_\_GPU\_METADATA} is a FlatBuffers document: a root table pointing at a per-kernel table with a vtable,
binding-record vectors, and further vectors and records (MM~\mm{25.143}, \code{agxforge/g17/mdgen.py}). Over 6,000 cached sections,
per-kernel slots 2, 4, 6, 8, 10, 12, 13, 26, 27 and 29 are pointers to length-prefixed vectors; 6, 8, 10 and 12 are empty
in 2,828 of 2,828; 2, 4 and 26 are built from the binding list (\code{agxforge/g17/authorobj.py}). Writing an integer into a
pointer slot creates a dangling pointer in a section the loader walks; that failure class hung this GPU on 2026–09-08,
and the author now refuses it. The per-kernel table declares 60 bytes, and every corpus object leaves at least 68
bytes before the next structure: packed to its 52 bytes of fields the loader segfaults, and a slack sweep put the
boundary between 8 and 12 bytes (\code{agxforge/g17/mdgen.py}). Binding records come in six measured shapes (long, short,
elided for buffer 0, written, mid, elided-written); across 9,207 records every slot set lies within \{0, 1, 2, 3\}, field 3
is the written flag, and no record carries a type: a buffer's element type never reaches \code{\_\_GPU\_METADATA}
(\code{agxforge/g17/mdgen.py}, \code{agxforge/g17/authorobj.py}).

The loader does not reach the binding records through the FlatBuffers root. In the measured scalar class the root's slot 3, the
only pointer that could reach the two binding vectors, is zero, yet the records are read: swapping two buffer indices
moved the store off its buffer and left the sentinel untouched. Whatever walks the section finds those records another
way (measured, \code{agxforge/g17/mdgen.py}). One field found inert on early kernels mattered on a later one: the
first two classes were built from blobs whose zeroing sweep had erased a whole single-record vector; that was harmless
while a kernel stored to a device buffer and fatal (a loader segfault) once a third binding record made slot 26 point at
the wrong vector (\code{agxforge/g17/mdgen.py}).

\Cref{tab:metadata-slots} lists the per-kernel slots whose meaning is established.

\begin{xltabular}{\linewidth}{@{}l>{\hsize=1.20\hsize}X>{\hsize=0.80\hsize}X@{}}
\caption{Per-kernel \code{\_\_GPU\_METADATA} slots with established meaning, with content and standing.}\label{tab:metadata-slots}\\
\toprule
Slot & Content & Standing \\
\midrule
\endfirsthead
\tablecontinued{3}
\toprule
Slot & Content & Standing \\
\midrule
\endhead
0 & declared register count = highest 32-bit register named in the main
program + 1 (2,022 of 2,022 corpus objects; 247 of 247 delivered
bundles) & Apple's compiler emits; inert for execution and
residency (\cref{sec:residency-occupancy-rules}; MM~\mm{10}, MM
\mm{25.115}); a corrupted slot-0 \emph{offset} crashes the driver (MM
\mm{25.147}) \\
1 & kind-6 field 2 + field 3 (\cref{sec:constant-program-record}) & Apple's compiler emits, law S1 (MM~\mm{25.120}) \\
2, 3, 4 & built from the bindings; in the scalar class slot 3 = 8 × bindings, slot
4 = Q - 4, slot 2 = Q - 4 + 4 × bindings, where Q is a per-kernel
multiple of 4 correlating with no readable field & measured: only the right slot-4 value executes
(\code{agxforge/g17/mdgen.py}) \\
13 & the constant pool, a vector whose count is a byte length (region 4 +
count + 16 bytes): loop trip counts, offsets, division-magic constants
(0xAAAAAAAB for /3, 0xCCCCCCCD for /5); changing one literal in the
source moves slot-13 bytes while \code{\_\_TEXT} is byte-identical & Apple's compiler emits (MM~\mm{25.118}, MM~\mm{25.120});
the uniform-halfword pool read by ordinary forms (
\cref{sec:constants-uniforms}) \\
15, 16, 17 & one byte each, value 1 whenever present; slot 16 = writes a device
buffer, 17 = writes a texture, 15 = 16 or 17: over 21,131 sections
exactly four combinations (none 286, texture only 70, buffer 20,757,
both 18). Threadgroup and imageblock stores set none & Apple's compiler emits (\code{agxforge/g17/cc.py}) \\
18 & threadgroup memory used (static or dynamic), value 1 when present (2,020
of 2,020) & Apple's compiler emits (\code{agxforge/g17/cc.py},
\code{agxforge/g17/mdgen.py}) \\
19 & the explicit imageblock declaration (with \code{\_\_GPU\_LD\_MD} slot 23
and an imageblock form of \code{\_\_GPU\_ARCH\_LD\_MD} carrying the
element size); in the tensor class slot 19 sits at 55 and slot 44 moves
to 54 & Apple's compiler emits (MM~\mm{25.89}, MM~\mm{25.104});
an undeclared image reads 0 (measured, 
\cref{sec:imageblock-per-core-tile}) \\
27 & binding-index words, a word count: 0, 1, 2 name buffers A, B, C & Apple's compiler emits (MM~\mm{25.118}, MM~\mm{25.120}) \\
28 & statically declared threadgroup bytes; Metal's
\code{staticThreadgroupMemoryLength} reports exactly this (0, 16, 64,
128, 16,896 over eight kernels) & measured (MM~\mm{25.147}) \\
29 & the system-register vector: 4 + 4n bytes, one entry per special register
read, ordered by entry value and keyed by register, not position; the
numbering is 
\cref{sec:slot-29-metadata-numbering}'s (for example
SR130, the lane index, is entry 52; SR156, threadgroup position x, is 0) & measured for the active classes (MM~\mm{0.8}, MM~\mm{15}) \\
31 & present iff the constant program writes at least one uniform (726 of 726
MMA programs; 22,299 of 22,305 objects); value 16 + 16 × ceil(buffers /
2) & Apple's compiler emits, laws L4 and S31 (MM~\mm{25.120}) \\
32 & instruction-count class: absent at 30 main-program instructions or
fewer, 3 with a back edge, otherwise 1 up to 300 and 2 from 301 (43
witnesses) & Apple's compiler emits (\code{agxforge/g17/mdgen.py}) \\
33 & present for programs with a runtime loop & Apple's compiler emits (MM~\mm{25.100},
\code{agxforge/g17/mdgen.py}) \\
38, 42 & the texture-state byte count, in a texture kernel & Apple's compiler emits (\code{agxforge/g17/mdgen.py}) \\
44 & the tensor-unit enable, one byte, value 1; present exactly when the code
contains an accelerator MMA (50 of 50 such kernels, no legacy-engine
kernel) & measured, causal over 14 pipelines (MM~\mm{6}; 
\cref{sec:ordering-enable}) \\
\bottomrule
\end{xltabular}

Clearing slot 44's low byte tests whether it is a validity flag, which the driver would check and refuse or fault on,
or an enable for the unit. The pipeline still builds and the dispatch completes without error, every tensor MMA
writes zeros, and the MMAs run fast: a 1,400-MMA chain takes 14.75~µs instead of 32.0 (8.4~ns per issue instead of 22), a
saturated dispatch 0.587~ms instead of 22.7. The MMAs still issue, at roughly ALU cost, but the unit does no work. Any
nonzero low byte, with the other three bytes arbitrary, is bit-identical to 1 (measured, MM~\mm{6}). The timing fits a
power or clock gate on the unit better than a validity check, though the mechanism was not measured
(\cref{sec:ordering-enable}). Since neither pipeline creation nor the dispatch rejects a cleared value, a missing enable
shows up only as zeros, and a compiler has to derive the slot from the presence of an accelerator MMA in the emitted code. Its
byte position in the per-kernel table depends on the class (an imageblock declaration moves it to 54, MM~\mm{25.89}; the
constant-program re-layout places it at 55, \cref{sec:constant-program-record}); "slot 44" in this document means the field, wherever the class puts it.

For slot 32, MM~\mm{10} calls the low byte inert, while MM~\mm{15} and \code{agxforge/g17/mdgen.py} treat slot 32 as an instruction-count property a body must satisfy;
the first concerns the value and the second its presence. The \emph{value's} low byte was mutated without
changing execution in the measured kernels (MM~\mm{10}). \emph{Presence} changes the table's layout: the tail's byte slots (33,
32, 16, 15) are right-aligned and the slots after them move, so a layout with the wrong presence places other fields,
among them slot 0 and slot 44, at the wrong offsets (\code{agxforge/g17/mdgen.py}, MM~\mm{25.120}). This compiler therefore emits
both presence and value by Apple's law (\cref{sec:program-boundaries} records the earlier padding to 31 instructions; MM~\mm{15}). Not established: whether a wrong value with the right presence matters for any kernel shape
beyond those tested.

\textbf{Worked example: the running example's contract and the metadata it implies} (compiled; delivered image executed
below Metal). Each row of \cref{tab:example-contract} gives a field of the contract this compiler records for the running example
(\code{example-listing.txt}, the line after the listing) and the metadata that field implies:

\begin{xltabular}{\linewidth}{@{}>{\hsize=0.69\hsize}X>{\hsize=0.83\hsize}X>{\hsize=1.48\hsize}X@{}}
\caption{Contract fields this compiler records for the running example, with the metadata each implies.}\label{tab:example-contract}\\
\toprule
Contract field & Value & Metadata consequence \\
\midrule
\endfirsthead
\tablecontinued{3}
\toprule
Contract field & Value & Metadata consequence \\
\midrule
\endhead
\code{register\_count} & 76 & slot 0 = 76: the highest register named is R75 (top of the B tile
R72\_R75) \\*
\code{main\_instruction\_count}, back edge & 101, none & slot 32 = 1 \\*
\code{system\_registers} & 130, 156 & slot 29 = {[}0, 52{]}: \op{14060} of
\code{SR\_SIMD\_ELEM} (the lane index) and \op{14059}
of \code{SR\_TG\_X} (threadgroup x) \\*
\code{execution.tensor} & true (four tensor MMAs) & slot 44 present, value 1 \\*
bindings & A (index 1, offset 0), B (2, offset 2), C (3, offset 4, written) & binding records, C in the written shape; descriptor bytes 0, 4, 8 (rank
x 4) in the code (MM~\mm{23.2}) \\*
\code{pk\_values} & slots 15 and 16 = 1 & writes a device buffer, no texture \\*
\code{constant\_pool}, \code{prologue} & empty; \code{0e000000} + thirty \code{0600} & no slot 31, no constant-program record; the 64 bytes before entry 64 are
the constant-program position, holding only \op{684}
and \op{13483} \\*
\code{launch} & exact grid required, not bounds-checked & runtime policy of the released class, not metadata \\*
\bottomrule
\end{xltabular}

\code{register\_count}, \code{main\_instruction\_count}, \code{forms}, the system-register vector and the launch facts come from the
emitted program and are never copied from a nearby kernel; \code{tensormetadata.layout} refuses an unwitnessed tail or
system-register set (this compiler, MM~\mm{23.2}). The contract's \code{register\_count} and \code{main\_instruction\_count} count
\code{\_agc.main} only (MM~\mm{25.146}).

Below Metal, the register count also travels in the launch state: the USC Registers packet (tag 0x8d) carries it in
bits 8–12 in groups of eight, plus a spill size (measured by capture; \cref{ch:submission-below-metal}).

\section{Constant program and record}\label{sec:constant-program-record}

Every Apple MMA program carries a second program, \code{agc.main.constant\_program} (with a \code{.cfg} suffix when
it branches), which runs once per dispatch, as one thread, before the main program, and computes thread-invariant values
into uniform registers that main reads (Apple's compiler emits; modelled and checked against Apple's GPU, MM~\mm{25.120},
MM~\mm{25.179}). Its inputs are binding entries read by \op{14061}: for the k-th buffer,
halfwords 8 + 4k and 10 + 4k hold its 64-bit pointer, low word then high word, in all 21 constant programs examined
(MM~\mm{25.59}, MM~\mm{25.179}; the two sections' notations agree, \cref{app:a}). \emph{Correction:} MM~\mm{25.113} read this metadata as a
register-spill description; it describes the constant program (MM~\mm{25.120}).

A constant-program instruction writes a uniform exactly when Apple's decoder
prints its first operand as \code{bin(op0, K, S)} (\opn{0} being the \op{0}); it writes uniform
K // 2 (the uniform file itself is \cref{sec:register-classes}). Writers are \op{592} in both forms and ALU forms
with a uniform destination, among them \op{10298}, \op{11185}, \op{1038} and \op{2318}; \cref{app:c} lists the rest (MM~\mm{25.120}). Three pair exactly with main-program forms
plus one extra mode operand: the integer-to-float writer with \op{11179}, the uniform / staging
write with \op{586}, and the fp32 add writer with \op{1006}
(MM~\mm{25.179}). K is 2 × ((byte0 >{}> 4) \textbar{} ((byte2 >{}> 6) \& 3) <{}< 4) in the 4-byte uniform / staging write, reaching 126 (\cref{sec:framing}), and a nine-bit field in the 8-byte and longer forms (MM~\mm{25.120}). The
worked bytes \code{8b040924}
write u8 (K 16) from R2 and wait on slot 0; \code{b308070256a0a402} writes u9 (K 18) from R0 and waits on slot 1 (decoder,
MM~\mm{25.120.1}).

\Cref{tab:constant-program-laws} gives the laws of the constant program's metadata (Apple's compiler emits; each checked on
every MMA program it applies to with zero misses, and predicted before compiling held-out variants, MM~\mm{25.120}).

\begin{xltabular}{\linewidth}{@{}lX@{}}
\caption{Laws relating the constant program to its metadata records and per-kernel slots.}\label{tab:constant-program-laws}\\
\toprule
Law & Statement \\
\midrule
\endfirsthead
\tablecontinued{2}
\toprule
Law & Statement \\
\midrule
\endhead
L4, S31 & per-kernel slot 31 present iff the constant program writes a uniform;
value 16 + 16 × ceil(buffers / 2), buffers = len(argument map) / 4 (32
or 48 in the corpus; 64 at 5 buffers, held out) \\*
L1 & kind-3 record field 0 = 8 + 4 × len(argument map) \\*
L2 & kind-3 field 2 = the constant program's register count
(1 if it writes no uniform) \\*
L5 & the argument map (count at record + 20) is {[}4b, 4b+1, 4b+2, 4b+3{]}
per buffer b whose binding entry the constant program reads \\*
L3 & kind-6 field 3 = 1 + max(K // 2) over every uniform write \\*
S1 & per-kernel slot 1 = kind-6 field 2 (the main program's
uniform count) + field 3 \\*
\bottomrule
\end{xltabular}

L3 was reached by refutation: "6 + the number of writes" failed on a variant that skipped u6-u7; "1 + the highest
uniform written by \op{592}" failed where a 32-bit add immediate into the uniform/staging file
wrote the top uniform; four compile rounds with committed predictions converged on "every writer counts" (MM~\mm{25.120}). Over the wider corpus L3 fits 2,650 of 2,720 programs with a constant-program write;
all 70 misses are non-MMA and under-predict, and in 50 of them the main program reads uniforms up to Apple's value, a
rule not pinned. (The u0-u5 reservation in MMA programs is \cref{sec:register-classes}.)

With a constant program the serializer re-lays the per-kernel tail: slot 44 moves to byte 55, slot 31 (a 32-bit
value) goes at 56, byte slots 33, 32, 16 and 15 are right-aligned to end at 63, and slot 0 goes at 64; the argument map
follows the kind-3 body; every later position moves (\code{relocate\_layout}), and Q is recomputed. All 726 Apple MMA programs
re-serialize byte-identically after the change (this compiler, MM~\mm{25.120}).

Of 726 Apple MMA programs, 265 have metadata sections this compiler reproduces byte for byte from recovered facts alone,
286 more reproduce only when Apple's slot-13 and slot-27 values are supplied, 72 differ and 103 are
refused as class limits (bindings beyond the class, unmeasured system-register sets, missing records) (MM~\mm{25.120}). The
section's size is fully determined by table shapes, the lengths of the per-kernel vectors (slots 13, 14, 26, 27, 29, 41)
and the constant-program record's length (56 classes, none ambiguous; MM~\mm{25.113}). The two blocking inputs, the slot-13
prefix values and the argument map, are facts about the kernel \emph{source}: one program reads 301/400/500 where its sibling
reads 300/400/500; only 631 of 1,824 nonzero prefix words appear anywhere in the code; two programs with byte-identical
argument reads list different buffers. They cannot be recovered from Apple's code, and for a program this compiler builds
they are known by construction (MM~\mm{25.120}).

The project's emulator, running the constant program first as one thread over the same buffers and handing
its uniform writes to main, equals Apple's GPU byte for byte on 6 of 6 sources with full-width random data, and every
deformed model (all published uniforms 0; the argument read naming the next binding; halfwords swapped) differs on all 6
(measured against Apple's GPU, MM~\mm{25.179}).

\textbf{Not known.} Whether a uniform write wider or narrower than 32 bits occupies exactly one uniform (the variants built to
test it emitted no such write); the main-program rule behind the 70 L3 misses; 442 constant-program uniform writes no
main program reads (MM~\mm{25.120}, MM~\mm{25.120.1}).

\section{ABI classes}\label{sec:abi-classes}

A class is a measured layout plus the free values inside it. This project authors only classes it has witnessed and
refuses the rest (this compiler, \code{agxforge/g17/mdgen.py}, \code{tensormetadata.layout}).

The released tensor class is ABI version 5 (this compiler's platform contract, MM~\mm{23.2}):

\begin{lstlisting}
ABI version                 5
entry                       64
execution                   {simd_width: 32, tensor: true}
public buffers              A=1, B=2, C=3
binding ranks in code       A=0, B=1, C=2  -> descriptor bytes 0, 4, 8
pointer-block offsets       0, 2, 4
system-register sets        (130,) or the measured composed class (130, 156)
constant pool               explicitly empty
pk_extra / pk_values        (15, 16) / both 1
prologue                    0e000000 followed by thirty 0600 fillers
grid / threadgroup          (32,1,1) / (32,1,1)
\end{lstlisting}

Three of its properties are machine rules (MM~\mm{23.2}). SR130 must be read in a tensor kernel. Descriptor byte values
are 4 × binding rank, where the rank is the position in the compiler's ordered binding contract and differs from the
public Metal buffer index. And without the tensor enable (slot 44) every MMA writes zeros (MM~\mm{6}). The single-simdgroup, exact-grid launch is runtime policy: research lowerings ran 2 and 4
simdgroups and up to 8 threadgroups bit-exactly (MM~\mm{0.13}, MM~\mm{25.89}).

Other classes have also been measured:

\begin{itemize}
\item The scalar object class (500-byte section, nine per-kernel slots, Q carried by class) and the tensor class (412 bytes,
  slot 44 at 55) are the two base layouts; further classes add binding shapes (the written-buffer class; reproduction counts are in
  \cref{sec:constant-program-record}), the composed system-register sets (130, 156), (130, 133, 156) and the imageblock set (130,
  164, 165), the constant-program re-layout, and loop tails without slot 32 (\code{agxforge/g17/mdgen.py}, MM~\mm{25.104}, MM~\mm{25.108},
  MM~\mm{25.120}).
\item The system-register class is keyed on the full set. \code{(SR130, SR156)} authors and runs as a measured class; it is not
  a universal rule for every form whose decoder byte 1 is 130 (13 of 385 such forms cannot inherit entry 52 from the
  coarse key) (MM~\mm{0.13}, MM~\mm{0.14}).
\item An empty system-register set must be refused: 7 real programs declaring one crashed the driver inside pipeline
  creation (measured, MM~\mm{25.147}).
\end{itemize}

The delivered manifest (\code{manifest.json}, format \code{g17-common-pipeline-v3}) records the ABI version, system registers,
per-binding index, offset, written flag and element type, the entry, and the (opcode, length) forms. 222 of 247 manifests
carry a template's \code{register\_count} and \code{code\_size}; the true values are in the container, and a tool reading the
manifest's count is misled (MM~\mm{25.147}).

The other metadata sections were measured by one-at-a-time mutation (\code{agxforge/g17/authorobj.py}):

\begin{itemize}
\item \code{\_\_GPU\_LD\_MD}: format known and entry PC named causally, but its slot set is a compiler input. A foreign but
  structurally valid \code{\_\_GPU\_LD\_MD} (a buffer kernel's in a texture kernel) builds a pipeline without complaint, which is
  why this section is never defaulted. It carries a per-process build counter that reads like a field (MM~\mm{25.91}).
\item \code{\_\_GPU\_ARCH\_LD\_MD}: one boolean, set in 6,234 of 21,001 objects; no contract fact and no opcode determines it (the
  best of seven tested predictors left 3,144 exceptions), so it is a compiler input.
\item \code{\_\_GPU\_STATS\_MD}: 96 bytes with 20,999 distinct patterns in 20,999 objects, not read by the loader: zeros, 0xFF,
  another kernel's bytes or removal all left an authored texture read exact.
\item Damaging \code{\_\_GPU\_METADATA}, \code{\_\_GPU\_LD\_MD} or \code{\_\_GPU\_ARCH\_LD\_MD} kills Apple's driver inside pipeline creation.
\end{itemize}

\section{Object, library and archive bytes}\label{sec:object-library-archive}

A bundle's \code{scan.arc.metallib} is a fat container holding an \code{air64} MetalLib and an \code{applegpu} Mach-O
GPU executable. It is loaded with \code{newBinaryArchiveWithDescriptor} and \code{newComputePipelineStateWithDescriptor} under
\code{MTLPipelineOptionFailOnBinaryArchiveMiss}, which makes Metal use the precompiled code rather than compile the AIR
(MM~\mm{25.143}). The project's tools author each layer: \code{g17mtlb.py} (the MetalLib), \code{g17container.py} (the Mach-O head and
AIR tables), \code{g17arc.py} (the archive's \code{\_\_reflection}, \code{\_\_compute}, \code{\_\_descriptor}, \code{\_\_metallib}), \code{g17authorobj.py}
(the object), \code{g17mdgen.py} (\code{\_\_GPU\_METADATA}), \code{g17ldmd.py} (\code{\_\_GPU\_LD\_MD}) (MM~\mm{25.143}). Metal caches pipelines by AIR
function hash, so a patched archive must be built in one process and dispatched in another, with a positive control;
without one, a "no change" result cannot distinguish an inert bit from a cached pipeline (measured, MM~\mm{17} rule 36).

A census (\code{tools/g17imagecensus.py}, MM~\mm{25.148}) assigns every byte of 421 distinct delivered archives (12,770,955
bytes, 1,253 index entries over ten delivery roots) to exactly one kind (\cref{tab:image-census}), and the tool refuses to
finish if a byte is left over.

\begin{xltabular}{\linewidth}{@{}>{\hsize=1.17\hsize}Xl>{\hsize=0.83\hsize}X@{}}
\caption{Bytes of 421 delivered archives by kind, with each kind's share and how it is checked (MM~\mm{25.148}).}\label{tab:image-census}\\
\toprule
Kind & Share & What is checked \\
\midrule
\endfirsthead
\tablecontinued{3}
\toprule
Kind & Share & What is checked \\
\midrule
\endhead
authored: \code{\_\_text} program bytes and prologue,
\code{\_\_GPU\_METADATA} & 85.6\% & program bytes equal the bundle's \code{program.bin} \\*
generated: fat header, both MetalLibs, container head, object header and
symbols, \code{\_\_GPU\_LD\_MD}, \code{\_\_GPU\_ARCH\_LD\_MD},
\code{\_\_text} filler & 10.6\% & the whole envelope is re-emitted and compared \\*
class: constants execution shows are read & 0.2\% & each compared at its offset (72 bytes) \\*
inert: constants nothing reads & 0.3\% & each is zero (76 bytes) \\*
tail: the descriptor's (392, 16) & 0.1\% & its value; inert, and not one of the 148 \\*
zeroed: \code{\_\_reflection}, \code{\_\_descriptor},
\code{\_\_GPU\_STATS\_MD}, \code{\_\_GPU\_REMARKS\_MD} & 3.2\% & all zero \\*
unexplained & 0 &  \\*
\bottomrule
\end{xltabular}

The envelope is a function of the object and the entry name: the module hash is sha256(object + name), both MetalLibs are
built from (name, hash), and the archive is re-emitted byte for byte in 421 of 421. \code{\_\_text} is prologue + \code{program.bin}
+ filler in 421 of 421. Rebuilding from recorded recipes reproduces 1,249 of 1,253 entries byte for byte; the other four
are stale roots whose recipe now compiles a different program, not authoring failures (MM~\mm{25.148}).

Apple's own reader (\code{xcrun metal-objdump --macho --private-headers}) names the load commands, which identifies most of the
148 class constants: what the project had called LC\_TABLE (0x31) is LC\_NOTE, whose \code{data\_owner} is the AIR name; 0x32 is
LC\_BUILD\_VERSION. With the TOOL\_* values from \code{<mach-o/loader.h>} most constants are toolchain identifiers
(\cref{tab:class-constants}; MM~\mm{25.148}).

\begin{xltabular}{\linewidth}{@{}llllX@{}}
\caption{Class constants by byte offset in slice 1 of the archive, with whether execution reads them, their value and
meaning (MM~\mm{25.148}).}\label{tab:class-constants}\\
\toprule
Offset & Bytes & Kind & Value & Meaning \\
\midrule
\endfirsthead
\tablecontinued{5}
\toprule
Offset & Bytes & Kind & Value & Meaning \\
\midrule
\endhead
824 & 4 & read & 1 & LC\_BUILD\_VERSION platform = PLATFORM\_MACOS \\
828, 832 & 8 & inert & 0 & minos, sdk (Apple writes 0x1A0600 = 26.6 as minos) \\
836 & 4 & read & 2 & ntools \\
840, 844 & 8 & inert & 0 & tools{[}0{]} (Apple writes TOOL\_GPUARCHIVER 0x407, version 0x7D17FF00 =
32023.255) \\
848, 852 & 8 & read & 0x403, 0x20001 & tools{[}1{]}: TOOL\_AIRNT\_PLUGIN, 2.0.1 \\
1088 & 32 & inert & 0 & the AIR\_DESCRIPTOR key (Apple's varies) \\
1260 & 32 & read & e3b0c442\dots{} & SHA-256 of the empty string, checked by value \\
1296, 1300 & 8 & read & 0x402, 0x403 & TOOL\_AIRNT, TOOL\_AIRNT\_PLUGIN \\
1304 & 4 & read & 0xC0000000 & open: required, meaning unknown \\
1324 & 4 & read & 1 & open: required, meaning unknown \\
1428 & 4 & read & 0x408 & TOOL\_METAL\_FRAMEWORK \\
1440 & 4 & read & 0x40000004 & open: required, meaning unknown \\
\bottomrule
\end{xltabular}

60 of the 72 read bytes are named; 12 are open. They belong to a record with a position rule, not to fixed slots: over
22,721 Apple archives the eight-word record starting with TOOL\_AIRNT sits at 1296 in 22,625, one word later in 58, and is
absent in 38, and the word at 1440 is missing in exactly the 96 archives where the record moved (MM~\mm{25.148}).
\emph{Correction:} an older "76 inert" list and an older one-kernel "image is ours" census are superseded by this census
(\cref{app:a}).

The GPU executes the machine code the project authored, from the copy Metal places in its heap. Two things were
tested: whether Metal recompiles the kernel from the archive's AIR (or serves a pipeline cached by AIR hash), and
whether a change to that heap copy after pipeline creation reaches the program that runs.

\begin{itemize}
\item \code{program.bin} appears byte for byte in the archive's \code{applegpu} slice. 1,620 operand-byte flips, each checked to
  leave the opcode stream unchanged on decode, were written into that region (no rebuild, AIR and its hash unchanged) and
  dispatched in a fresh process; all 2,048 output floats and the kernel's time changed (MM~\mm{25.143}). The executed code
  is the archive's machine code, not a recompilation of its AIR.
\item For 10 delivered kernels of six kinds, the program's exact bytes appear at offset +0x2740 of the dispatch's first
  mapped 64 KiB region. The causal test patched that heap copy after \code{newComputePipelineState} and before the dispatch,
  on \code{C[i] = 3*i + 1}, changing the add constant (program byte 19) from 1 to 2 only if the bytes occurred exactly once.
  Unlike the operand flips above, this flip is in a byte adjacent to an opcode.
  The unpatched control read \code{3i + 1}; the heap patch read \code{3i + 2}, the same result as the flip authored in the
  container; and a patch asking for absent bytes wrote nothing and read \code{3i + 1}, so the patcher alters nothing else
  (measured, MM~\mm{25.147}). This establishes that the executed program incorporates the patched bytes; it does not determine
  whether subsequent copying or validation occurs.
\end{itemize}

Metal reports \code{maxTotalThreadsPerThreadgroup} 1,024 and \code{threadExecutionWidth} 32 for every kernel and
\code{staticThreadgroupMemoryLength} equal to slot 28; the declared register count it reports is the container's own
(measured, MM~\mm{25.147}). Where the driver takes the thread limit from is \cref{sec:threads-simdgroups-threadgroups} (MM~\mm{25.141.20}).

What pipeline creation validates was mapped by handing it the known-answer container with one change at a time
(\code{spike/agxsub/g17validprobe.py}; nothing dispatched; MM~\mm{25.147}). Two of the inputs are negative controls that
pipeline creation should reject. The outcomes:

\begin{itemize}
\item \emph{Rejected} ("invalid format"): a corrupted archive header, an archive truncated to half. These controls show that
  pipeline creation does check its input, so the acceptances below are not explained by an absent check.
\item \emph{Accepted}: code replaced by undecodable \code{0xff}, the end of program removed, declared register count 0, 1, 200 or 256,
  static threadgroup memory declared at 1 MiB (reported back as 1,048,576, far past the 32 KiB a threadgroup can have).
\item \emph{Crashed the driver} (a CPU fault inside \code{newComputePipelineState}): a corrupted per-kernel slot-0 offset, and the 7
  programs with an empty system-register set.
\end{itemize}

Since pipeline creation checks neither the instruction bytes nor the metadata values, and a bad metadata layout faults
the driver, errors inside the format have to be caught when the image is authored. This project's guards (
\cref{sec:admission-decodes-executes}) and refusals run before the image is written for that reason.

Code generation and the whole container are this project's; no Apple compiler process
runs to produce them, and Apple's decoder is read in-process only as a check. Turning Metal source into AIR still runs
Apple's front end, and the knowledge of which forms exist was learned from Apple's compiler output (MM~\mm{25.147}).

\textbf{Evidence.} MM~\mm{0.8}, MM~\mm{0.13}, MM~\mm{0.14}, MM~\mm{6}, MM~\mm{10}, MM~\mm{15}, MM~\mm{17} (rule 36), 23.2, 23.5, 25.59, 25.89, 25.91, 25.100, 25.104, 25.108,
25.113, 25.118, 25.120, 25.120.1, 25.141.20, 25.143, 25.146, 25.147, 25.148, 25.179; \code{agxforge/g17/mdgen.py},
\code{agxforge/g17/authorobj.py}, \code{agxforge/g17/cc.py}.
